\clearpage\corrigesChapitre{18}{Lentilles et vision}
\corrigesSouples
\begin{blocLivre}{cor-18-1}\begin{corrige}{1}{Repérer les foyers}
\schema[corrige]{lentille-foyers}
\begin{questions}
\item À 10 cm de part et d’autre de O, sur l’axe optique.
\item Le foyer image $F^{\prime}$, du côté de sortie de la lumière.
\end{questions}
\end{corrige}\end{blocLivre}
\begin{blocLivre}{cor-18-2}\begin{corrige}{2}{Le rayon central}
\begin{questions}
\item On le prolonge en ligne droite, sans déviation.
\item Non : il faut au moins une autre direction de rayon pour déterminer leur intersection.
\end{questions}
\end{corrige}\end{blocLivre}
\begin{blocLivre}{cor-18-3}\begin{corrige}{3}{Un rayon parallèle}
\begin{questions}
\item Par le foyer image $F^{\prime}$.
\item Un rayon incident qui passe par le foyer objet F ressort parallèle à l’axe.
\end{questions}
\end{corrige}\end{blocLivre}
\begin{blocLivre}{cor-18-4}\begin{corrige}{4}{Lire une image réelle}
\schema[corrige]{lentille-2f}
\begin{questions}
\item Le rayon central et le rayon parallèle réfracté par $F^{\prime}$ se coupent à droite ; la version corrigée montre la construction.
\item L’image est à 20 cm de O, de taille 2,0 cm, renversée.
\end{questions}
\end{corrige}\end{blocLivre}
\begin{blocLivre}{cor-18-5}\begin{corrige}{5}{Un grandissement}
\begin{questions}
\item $\gamma=-1,5/3,0=-0,50$.
\item Elle est réduite car $|\gamma|=0,50<1$.
\end{questions}
\end{corrige}\end{blocLivre}
\begin{blocLivre}{cor-18-6}\begin{corrige}{6}{Construire une image réduite}
\schema[corrige]{lentille-3f}
\begin{questions}
\item Les rayons se coupent à 15 cm après la lentille.
\item L’image mesure 1,0 cm et est renversée : $\gamma=-1,0/2,0=-0,50$.
\end{questions}
\end{corrige}\end{blocLivre}
\begin{blocLivre}{cor-18-7}\begin{corrige}{7}{Le modèle de l’œil}
\schema[corrige]{oeil}
\begin{questions}
\item Pupille : diaphragme ; système cornée-cristallin : lentille ; rétine : écran.
\item Dans le plan de la rétine.
\end{questions}
\end{corrige}\end{blocLivre}
\begin{blocLivre}{cor-18-8}\begin{corrige}{8}{Utiliser Thalès}
\begin{questions}
\item $|\gamma|=p^{\prime}/p=24/12=2,0$.
\item La taille vaut 2,0 cm et l’image réelle d’un objet réel dans ce montage est renversée.
\end{questions}
\end{corrige}\end{blocLivre}
\begin{blocLivre}{cor-18-9}\begin{corrige}{9}{Un écran mal placé}
\begin{questions}
\item Les rayons d’un même point ne se rejoignent pas dans le plan actuel de l’écran.
\item Oui : rapprocher l’écran de la lentille jusqu’au plan de l’image, sans modifier les autres positions.
\end{questions}
\end{corrige}\end{blocLivre}
\begin{blocLivre}{cor-18-10}\begin{corrige}{10}{Cacher une moitié}
\begin{questions}
\item Non. Chaque point de l’image peut encore recevoir des rayons passant par la partie libre de la lentille, si elle les laisse passer.
\item L’image devient moins lumineuse car moins de lumière est transmise.
\end{questions}
\end{corrige}\end{blocLivre}
\begin{blocLivre}{cor-18-11}\begin{corrige}{11}{Une accommodation}
\begin{questions}
\item Il modifie sa convergence par accommodation, principalement grâce au cristallin.
\item Non : l’analogie avec le banc aide à comprendre, mais l’œil ne déplace pas sa rétine pour faire la mise au point.
\end{questions}
\end{corrige}\end{blocLivre}
\begin{blocLivre}{cor-18-12}\begin{corrige}{12}{Valider un tracé}
\demarche{Former le quotient distance image sur distance objet et traiter le sens séparément.}
\begin{questions}
\item $p^{\prime}/p=15/30=0,50$, et non 2. Les distances ont été inversées dans le quotient.
\item $\gamma=-0,50$. Le signe traduit le renversement.
\end{questions}
\retiens{Un grandissement négatif correspond ici à une image renversée.}
\end{corrige}\end{blocLivre}

